\fancyhead[L]{\small Beiträge zur Theorie}
\SourceLeaf{080}
\begin{center}
{\large\bfseries BEITRÄGE ZUR THEORIE\par}
\medskip
{\Large\bfseries DER ALGEBRAISCHEN GLEICHUNGEN.\par}
\medskip
\rule{0.18\linewidth}{0.4pt}
\end{center}

In dieser Abhandlung werden zwei verschiedene Gegenstände behandelt, die algebraische Gleichungen betreffen. Zuerst stelle ich den vor fünfzig Jahren von mir gegebenen Beweis des Fundamentalsatzes der Theorie algebraischer Gleichungen in veränderter Form und mit erheblichen Zusätzen dar. Der zweite Teil ist einer besonderen Behandlung algebraischer Gleichungen mit drei Gliedern gewidmet und enthält Methoden, mit denen sich nicht nur die reellen, sondern auch die imaginären Wurzeln solcher Gleichungen leicht bestimmen lassen.

\begin{center}
\rule{0.18\linewidth}{0.4pt}\par
\medskip
\textbf{ERSTE ABTEILUNG.}
\end{center}

Die im Jahr 1799 erschienene Abhandlung \emph{Ein neuer Beweis des Satzes, dass jedes Polynom in einer Variablen in reelle Faktoren ersten oder zweiten Grades zerlegt werden kann} hatte zwei Ziele: erstens zu zeigen, dass sämtliche bis dahin versuchten Beweise dieses wichtigsten Satzes der Theorie algebraischer Gleichungen unzureichend und illusorisch sind, und zweitens einen neuen, vollkommen strengen Beweis zu geben. Es ist nicht nötig, auf den ersten Gegenstand noch einmal zurückzukommen. Dem dort gegebenen neuen Beweis habe ich selbst später noch zwei andere folgen lassen; ein vierter wurde zuerst von \textsc{Cauchy} aufgestellt. Diese vier Beweise \SourceLeaf{081}beruhen alle auf ebenso vielen verschiedenen Grundlagen. Sie stimmen aber darin überein, dass jeder zunächst nur das Vorhandensein \emph{eines} Faktors der betreffenden Funktion beweist. Der Strenge der Beweise tut dies allerdings keinen Abbruch. Denn es ist klar, dass nach dem Abtrennen dieses einen Faktors von der gegebenen Funktion eine ähnliche Funktion niedrigeren Grades übrig bleibt, auf die der Satz erneut angewandt werden kann, und dass die Wiederholung des Verfahrens schließlich eine vollständige Zerlegung der ursprünglichen Funktion in Faktoren der bezeichneten Art ergeben muss. Dennoch gewinnt jeder Beweis zweifellos an Vollendung, wenn gezeigt wird, dass er das Vorhandensein sämtlicher Faktoren unmittelbar anschaulich machen kann. Dass dies beim ersten Beweis der Fall ist, habe ich bereits in der genannten Abhandlung, Artikel 23, angedeutet, ohne es dort weiter auszuführen. Dies soll jetzt ergänzt werden. Zugleich nutze ich die Gelegenheit, die Hauptpunkte des gesamten Beweises in veränderter und, wie ich glaube, größerer Klarheit bietender Form zu wiederholen. Was die äußere Formulierung des Satzes selbst betrifft, so wurde 1799 die Aussage, dass die Funktion $x^n+Ax^{n-1}+Bx^{n-2}+\mathrm{usw.}$ in reelle Faktoren ersten oder zweiten Grades zerlegt werden kann, deshalb gewählt, weil jede Einbeziehung imaginärer Größen vermieden werden sollte. Heute, da der Begriff der komplexen Größen jedem geläufig ist, scheint es angemessener, jene Form aufzugeben und den Satz so auszusprechen, dass die Funktion in $n$ \emph{lineare} Faktoren zerlegt werden kann. Die konstanten Teile dieser Faktoren brauchen dann keine reellen Größen zu sein; für sie müssen auch beliebige komplexe Werte zugelassen werden. Bei dieser Formulierung gewinnt der Satz selbst an Allgemeinheit, weil auch bei den Koeffizienten $A,B$ usw. keine Beschränkung auf reelle Werte vorausgesetzt werden muss, sondern beliebige Werte zulässig bleiben.

\FifthArticle{1}
Wir betrachten demnach die Funktion der Unbestimmten $x$
\[
x^n+Ax^{n-1}+Bx^{n-2}+\mathrm{usw.}+Mx+N=X
\]
wobei $A,B\,.\,.\,.\,.\,M,N$ bestimmte reelle oder imaginäre Koeffizienten bezeichnen.
\par
\SourceLeaf{082}
Aus der elementaren Algebra ist der Zusammenhang zwischen den Wurzeln der Gleichung $X=0$ und den linearen Faktoren von $X$ bekannt. Wird nämlich jene Gleichung durch die Substitution $x=p$ erfüllt, so ist $x-p$ ein Faktor von $X$. Gibt es $n$ verschiedene Möglichkeiten, jene Gleichung zu erfüllen, nämlich $x=p$, $x=p'$, $x=p''$ usw., so ist das Produkt $(x-p)$ $(x-p')(x-p'')\, .\, .\, .\, .$ mit $X$ identisch. Unter besonderen Umständen kann aber eine Lösung wie $x=p$ in $X$ auch den Faktor $(x-p)^2$ oder $(x-p)^3$ oder eine höhere Potenz bedingen; dann zählt man die Wurzel $p$ als zweimal, dreimal usw. vorhanden.

Soll also nur bewiesen werden, dass die Funktion $X$ sicher \textit{einen} linearen Faktor hat, so genügt es, das Vorhandensein irgendeiner Wurzel der Gleichung $X=0$ nachzuweisen. Soll dagegen die vollständige Zerlegbarkeit der Funktion in lineare Faktoren auf einmal bewiesen werden, so muss gezeigt werden, dass die Gleichung $X=0$ entweder durch $n$ ungleiche Werte von $x$ oder durch eine kleinere Anzahl ungleicher Lösungen erfüllt werden kann, von denen aber ein Teil die Eigenschaften mehrfach zählender gleicher Wurzeln so besitzt, dass die Zusammenzählung aller ungleichen und gleichen Wurzeln die Gesamtsumme $=n$ ergibt.

